\section{引言}

随着智能交通系统的快速发展，车辆轨迹数据的精确获取与优化成为提升交通管理与安全的关键环节。卡尔曼滤波作为一种经典的递归估计算法，广泛应用于动态系统的噪声抑制与状态重构。本文面向真实城市场景，针对车辆轨迹观测中常见的传感器噪声与局部遮挡问题，提出一种基于自适应卡尔曼滤波的轨迹优化框架。

研究构建了一个包含位置、速度与加速度三阶运动模型的状态空间方程，并通过实时噪声协方差估计增强滤波器鲁棒性。在公开数据集KITTI与自采集城市交叉口数据上的对比实验表明，所提方法相较于标准卡尔曼滤波与扩展卡尔曼滤波，在均方根误差（RMSE）与轨迹平滑性两项指标上分别降低了约15%与20%（详见第四节）。本文主要贡献可概括为以下三点：其一，设计了可在线调整系统噪声协方差的机制，克服了固定噪声参数在动态环境下的退化问题；其二，引入观测异常检测模块，有效滤除由于遮挡引起的突发野值点；其三，结合车辆运动学约束的预测-更新循环，实现了端到端的离线与在线轨迹优化方案。

本文其余部分安排如下：第二节介绍问题建模与自适应滤波器的详细推导；第三节描述实验设计方案及数据集准备；第四节展示定量与定性分析结果；第五节给出结论并对未来工作进行展望。